% ----------------------------------------------------------------------------
%  Page de garde. Mise en page de l'en-tête institutionnel INSAT, adaptée à un
%  projet académique : ni entreprise, ni encadrant, ni page de validation.
% ----------------------------------------------------------------------------
\begin{titlepage}
\thispagestyle{empty}
\centering
\singlespacing

\begin{minipage}[c]{0.2\textwidth}
  \includegraphics[width=2.8cm]{logo_insat.jpg}
\end{minipage}\hfill
\begin{minipage}[c]{0.56\textwidth}
  \centering\footnotesize
  \textsc{République Tunisienne}\\
  Ministère de l'Enseignement Supérieur et de la Recherche Scientifique\\[2pt]
  Université de Carthage\\[2pt]
  \textbf{Institut National des Sciences Appliquées et de Technologie}
\end{minipage}\hfill
\begin{minipage}[c]{0.2\textwidth}
  \raggedleft\includegraphics[width=2.6cm]{logo_ucar.png}
\end{minipage}

\vspace{2pt}\rule{\textwidth}{0.6pt}

\vspace{0.9cm}
{\LARGE\textsc{Rapport de Projet}\par}
\vspace{0.5cm}
{\large Sciences Appliquées et Technologie\par}
\vspace{0.3cm}
{\large Spécialité : Réseaux Informatiques et Télécommunications\par}
{\large Niveau : 5\textsuperscript{e} année (RT5)\par}

\vspace{0.7cm}
\rule{0.85\textwidth}{0.8pt}\\[0.35cm]
{\Large\bfseries\color{insatblue}
Communication sémantique apprise pour la transmission d'images et de vidéo sur canal sans fil simulé\par}
\vspace{0.25cm}\rule{0.85\textwidth}{0.8pt}

\vspace{0.8cm}
{\large Réalisé par\par}
\vspace{0.35cm}
{\large\bfseries
\begin{tabular}{c@{\hspace{2.2cm}}c}
EL DABBABI Mouine & SADDOUD Mohamed Amine \\[0.35cm]
BOUCHAALA Oussema & LAKHADER Mohamed Helmi \\
\end{tabular}\par}

\vfill
\begin{tabular}{@{}r@{\hspace{0.6em}}l@{}}
  \textbf{Projet :} & Projet 07 -- Codage conjoint source-canal appris (deep JSCC)\\
  \textbf{Canaux étudiés :} & AWGN et Rayleigh par blocs, simulés\\
  \textbf{Jeux de données :} & CIFAR-10, Kodak, séquences vidéo CIF\\
\end{tabular}

\vspace{0.6cm}
{\large Année universitaire : \textbf{2025/2026}\par}
\end{titlepage}
